
\section{Exact seven-ray cone completion and the compulsory carrier}
\label{sec:seven-ray-cone-XXV}
We now test the smooth two-frequency covariance of Edition XXIV
against the actual H8 seven-ray geometry. These are exact
source-independent tensor identities in a local co-rotating frame,
not an unforced Navier--Stokes construction.

\begin{theorem}[Direct seven-ray cone obstruction]
\label{thm:ray-direct-XXV}
Let $t=1/4$, and in the local frame
$(e_1,e_2,e_3)=(e_r,e_\theta,e_z)$ put
\[
v_1=e_1,\qquad v_{2,3}=e_1\pm te_2,\qquad
v_{4,5}=e_1\pm te_3,\qquad
v_{6,7}=e_1+t(e_2\pm e_3).
\]
For $d_j\ge\mu>0$, the stress
$S=\sum_{j=1}^7d_jv_j\otimes v_j$ obeys
\[
\boxed{S_{rr}-S_{\theta\theta}\ge 105\mu/16>0.}
\]
On the Edition XXIV annular plateau, its manufactured two-frequency
Reynolds covariance has
$Q_{rr}-Q_{\theta\theta}=-2\lambda^2r^4\cos^2z\le0$.
Hence $Q+\pi I$ is not directly in the strict seven-ray cone for
any scalar gauge $\pi$.
\end{theorem}
\begin{proof}
Each ray satisfies $(v_j)_r=1$ and
$|(v_j)_\theta|\le1/4$. Consequently
$S_{rr}-S_{\theta\theta}
=\sum_jd_j(1-(v_j)_\theta^2)
\ge(15/16)\sum_jd_j\ge105\mu/16$.
The manufactured $Q$ has $Q_{rr}=2\cos^2z$ and
$Q_{\theta\theta}=2\cos^2z(1+\lambda^2r^4)$.
An isotropic gauge does not change their difference.
\end{proof}

\begin{theorem}[Strict seven-ray completion with a positive carrier]
\label{thm:ray-carrier-XXV}
On $3/4\le r\le5/4$, $|z|\le\pi/6$ and
$|\lambda|\le1/16$, write $s=\sin z$ and $c=\cos z$.
Define
\begin{align*}
d_1&=1,&
d_2&=1+4\lambda r^2c^2,&
d_3&=3-4\lambda r^2c^2,\\
d_4&=2+16s^2+16\lambda^2(r^2s^2-r^4c^2)
  +16\lambda rcs,\\
d_5&=2+16s^2+16\lambda^2(r^2s^2-r^4c^2)
  -16\lambda rcs,\\
d_6&=1-16\lambda rcs,&
d_7&=1+16\lambda rcs,\\
\pi&=\frac38-2c^2(1+\lambda^2r^4),\\
K&=\frac{85}{8}+32s^2+32\lambda^2r^2s^2
                     -30\lambda^2r^4c^2.
\end{align*}
For the complete co-rotating covariance tensor $Q$ of the
two-frequency field in Edition XXIV, the following holds exactly:
\begin{equation}\label{eq:ray-completion-XXV}
\boxed{Q+K e_r\otimes e_r+\pi I
=\sum_{j=1}^7 d_jv_j\otimes v_j,\qquad
d_j\ge\frac38>0,\quad K>0.}
\end{equation}
\end{theorem}
\begin{proof}
The covariance in the plateau is
\[
Q=2\begin{pmatrix}
c^2&\lambda r^2c^2&0\\
\lambda r^2c^2&c^2(1+\lambda^2r^4)&-\lambda rcs\\
0&-\lambda rcs&1+\lambda^2r^2s^2
\end{pmatrix}.
\]
The seven-ray synthesis identities are
\begin{align*}
S_{rr}&=\sum_jd_j,&
S_{\theta\theta}&=(d_2+d_3+d_6+d_7)/16,\\
S_{zz}&=(d_4+d_5+d_6+d_7)/16,&
S_{r\theta}&=(d_2-d_3+d_6+d_7)/4,\\
S_{rz}&=(d_4-d_5+d_6-d_7)/4,&
S_{\theta z}&=(d_6-d_7)/16.
\end{align*}
Substitution verifies all six symmetric tensor components.
Using $r\le5/4$, $|s|\le1/2$, $|c|\le1$ and
$|\lambda|\le1/16$ yields
\[
d_1=1,\quad d_2\ge39/64,\quad d_3\ge167/64,\quad
d_{4,5}\ge5007/4096,\quad d_{6,7}\ge3/8.
\]
Also $K\ge85/8-30(1/256)(625/256)>0$.
\end{proof}

\begin{proposition}[Non-gauge carrier curl]
\label{prop:ray-carrier-curl-XXV}
The spatially varying radial carrier in
\eqref{eq:ray-completion-XXV} has
\[
\nabla\cdot(K e_r\otimes e_r)=(K_r+K/r)e_r,
\]
and
\[
\boxed{\bigl[\nabla\times\nabla\cdot
(K e_r\otimes e_r)\bigr]_\theta
=\left(\frac{64}{r}+192\lambda^2r+300\lambda^2r^3\right)
\sin z\cos z.}
\]
The curl is nonzero on open subpatches with $0<|z|<\pi/6$.
Thus this carrier is not a pressure-gauge adjustment.
\end{proposition}
\begin{proof}
The cylindrical divergence of an axisymmetric pure radial
tensor gives $(K_r+K/r)e_r$.
Take its axial derivative and use the displayed formula for $K$.
A gradient has zero curl.
\end{proof}

\paragraph{Physical and independent-verification limits.}
This is an exact local positive cone completion, not the full
H8 four-field stress or a source-matched unforced PDE solution.
A positive radial carrier has a physical projected-divergence
effect and its construction must be justified; a pressure gauge
alone cannot supply it. The co-rotating ray frame also generates
spatial derivative terms. Source A.11/A.19, collar matching,
the seven-ray virtual-work dual, time evolution, all-mode
convergence and singularity inheritance remain unproved.
Independent symbolic reconstruction of all six tensor entries,
the rank-six tensor map with coefficient kernel
$(2,-2,0,-1,-1,1,1)$, five thousand exact rational parameter
replays, and an independent carrier-curl derivative audit pass.
The accompanying Lean declarations cover only finite scalar
identities, not the full PDE theorem.
